% 使用 ExBook 文档类，并传递选项
\documentclass[standard]{ExBook} 
 
\begin{document}

% 加载配置  
% 封面设置
\CoverImg{img/cov01.jpg} % 封面图片
\PreTitle{ExBook · 刷题本模板} % 前置标题
\Title{此处填写主标题} % 主标题
\TitleDescription{此处填写副标题} % 副标题
\TypeOne{A4紧凑版} % A4紧凑版下的类型标识
\TypeTwo{A4标准版} % A4标准版下的类型标识
\TypeThree{横版Pad版} % 横版Pad版下的类型标识
\TypeFour{A4宽松版} % A4宽松版下的类型标识
\TypeFive{A4单题版} % A4单题版下的类型标识
\TypeSix{竖版Pad版} % 竖版Pad版下的类型标识
\motto{你这个年龄是怎么睡得着觉的} % 封面座右铭
\Creator{研小布} % 制作人
\UpdateTime{\today} % 更新时间

% 页眉页脚设置
\Lhead{微信公众号·研小布} % 左页眉 
\Chead{2025考研} % 中页眉、平板模式（padl或padp）下页眉中间的文字
\Rhead{408WD数据结构选择题刷题本} % 右页眉、平板模式（padl或padp）下页眉右侧的文字
\LheadC{公众号·研小布·} % 平板模式（padl或padp）下页眉左侧的文字

% 水印设置
\TextWater{【微信公众号·研小布】} % 文字水印 出现在题目后面
\WaterImg{img/water.png} % 图片水印 出现在页面的右下角

% 主题颜色设置，可选主题有：
% blue (默认)
% green
% purple
% orange
% infj
% enfp
% infp
% esfp
% intj
% entp
% isfj
% enfj
\setThemeColor{\orange}

% 加载封面
\maketitle 
 
% 加载声明
\clearpage
\vspace*{0mm} % 添加两厘米的垂直空白
\begin{center}
    \large \heiti 前言
\end{center}

% \hspace*{2em}此刷题本依据 \underline{\textbf{\preBookName}}重新排版制作而成! 


\hspace*{2em}此刷题本的目的是方便大家在考研备考中多次刷题、记录自己的刷题过程和笔迹，以便日后复盘与巩固！
此刷题本不包含答案，答案请参考原书！

\hspace*{2em}关于其他考研刷题本的更新，请关注 \faWeixin 微信公众号：\underline{\textbf{研小布}}，可使用微信扫描下面的二维码关注：


\begin{minipage}[t]{1.0\textwidth}
    \centering
    \includegraphics[width=0.32\textwidth]{img/qr01.jpg} 
\end{minipage}

\hspace*{2em}如果你有\textbf{25考研数学}或者\textbf{专业课}清晰的PDF，可以投稿给小布，我来帮助制作刷题本，投稿微信：\textbf{\texttt{shang\_an\_001}}

\hspace*{2em}最后，希望此刷题本能够让你更高效地刷题，上岸值 ++ ！祝考研顺利。
\clearpage 

% 加载广告
\clearpage
\hideheaderfooter

\vspace*{-5mm} 
\begin{center}
    \large \heiti 打印纸质版说明\footnote[1]{此说明只针对A4版做题本（A4标准版、A4宽松版、A4紧凑版、A4单题版）}
\end{center}

\hspace*{2em}\heiti{打印参数建议~~} \songti{A4纸张、黑白（或彩色）、双面（或单面） }

\hspace*{2em}\heiti{打印渠道推荐~~} \songti{微信扫描下方二维码进入小程序可进行在线打印，超优惠打印价格！70gA4纸单面0.07元/张，
双面0.05元/张。}

\vspace*{4em}

\imgin[0.28]{}{img/print.png}

\clearpage

\setcounter{page}{1}
\tableofcontents 
    
\clearpage 

\section{文字录入型做题本第一章}
\subsection{文字录入型做题本第一章第1节}\qanswerloc{10}

\begin{qitems}

    \begin{bbox}
        \qitem 设$f(x)$满足$2f(x)+f(1-x)=x^2 abc$，则$ \displaystyle\sum \limits_{n}^{0}f(x)=\blankline.$
    \end{bbox}

    \begin{bbox}
        \qitem   设$f(x)=2x+\sqrt{x^{2}+2x+1}$，$g(x)=
        \begin{cases}
        x+2, & x\geqslant0, \\
        x-1, & x<0,  
        \end{cases}$，则$g[f(x)]= \blankline$.
    \end{bbox}

    \begin{bbox}
        \qitem   设某项目用于 发和宣传 总成本为$a$万元 当 发和宣传所 成本分别为$x$万元和 $y$ 万元时， 收益为$R=2x^{\frac{1}{3}}y^{\frac{1}{2}}$万元,则收 最大时，研发所用成本为\blankline.
    \end{bbox}

    \begin{bbox}
        \qitem   已知 $\lim\limits_{x\to0}\dfrac{f(x)}{x}$ 存在，且函数
        $$f(x)=\ln(1+x)+2x\bullet\lim_{x\to0}\frac{f(x)}{\sin x}$$
        则$\lim\limits_{x\to0}\dfrac{f(x)}{x}=$ \blankline.
    \end{bbox}

    \begin{bbox}
        \qitem  可以用\blankbox 定义一个完整的数据结构。\textwater
        \fourchoices{数据元素}{数据对象}{数据关系}{抽象数据类型}     
    \end{bbox}

    \begin{bbox}
        \qitem   若某算法的空间复杂度为$O(1)$，则表示该算法\blankbox 。
        \fourchoices
        {不需要任何辅助空间}
        {所需辅助空间大小与问题规模$n$无关}
        {不需要任何空间}
        {所需空间大小与问题规模$n$无关}
    \end{bbox}
\end{qitems}

\subsection{文字录入型做题本第一章第2节}
\qanswerloc{15}

\begin{qitems}

    \begin{bbox}
        \qitem   下列关于时间复杂度的函数中，时间复杂度最小的是\blankbox 。
        \fourchoices
        {$T_1(n)=n\log_2n +5000n$}
        {$T_2(n)=n^2 - 800n$}
        {$T_3(n)=n\log_2n - 6000n$}
        {$T_4(n)=20000\log_2n$}
    \end{bbox}

    \begin{bbox}
        \qitem   【2017 统考真题】 下列函数的时间复杂度是 \blankbox 。
        \begin{lstlisting}
    int func(int n){
        int i=0, sum=0;
        while(sum<n) sum += ++i;
        return i; 
    }
        \end{lstlisting}
        \fourchoices{$O(\log n)$}{$O(n^{\frac{1}{2}})$}{$O(n)$}{$O(n\log n)$}
    \end{bbox}

\end{qitems}

\section{文字录入型做题本第二章（长标题测试长标题测试）}
\subsection{文字录入型做题本第二章第1节（长标题测试长标题测试）}
\qanswerloc{20}

\begin{qitems}
    
    \begin{bbox}
        \qitem  已知曲线$L:y=\ln\sqrt{x}(2\leqslant x\leqslant4)$，在$L$ 上的任意点$P(x,y)$作切线，记切线与曲线$L$在 $2\leqslant x\leqslant4$
        时所围成的有界区域的面积为$S$.
        \begin{subqitems}
            \subqitem 求一点$P_0$，使上述面积$S$关于$x$的变化率为零；
            \subqitem 当点$P(x,y)$在曲线上移动至$(\mathrm{e},\dfrac{1}{2})$时，横坐标关于时间的变化率为1，求此时面积关于时间的变化率$\dfrac{\mathrm{d}S}{\mathrm{d}t}.$
        \end{subqitems}
    \end{bbox}

    \begin{bbox}
        \qitem 以 $yOz$ 面上的平面曲线段$y=f(z)(z\geqslant0)$ 绕$z$轴旋转一周所成旋转曲面与xOy 面围成一个无上盖容器(见图)，现以 3 cm$^3/$s 的速率把水注人容器内，水面的面积以$\pi$ c$m^2$/ s 的速率增大.已知容器底面积为 16$\pi$ c$m^2$,求曲线$y=f(z)$的方程.
        \imgin[0.4]{}{fig/fig01.png}
    \end{bbox}

    \begin{bbox}
        \qitem   分析以下各程序段， 求出算法的时间复杂度.
        \begin{lstlisting}[escapeinside={(*@}{@*)}]
    (*@\ding{172}:@*)
    i=1; k=0;
    while(i<n-1){
        k=k+10*i;
        i++;
    }

    (*@\ding{173}:@*)
    y=0;
    while((y+1)*(y+1)<=n)
    y=y+1;

    (*@\ding{174}:@*)
    for(i=0;i<n;i++)
        for(j=0;j<m;j++)
            a[i][j]=0;
    
        
        \end{lstlisting}
    \end{bbox}

    \begin{bbox}
        \qitem   【2011 统考真题】一个长度为 $L$（$L\geqslant 1$ ）的升序序列$ S$, 处在第$\lceil L/2\rceil $个位置的数称为 $S$
        的中位数。例如，若序列 $S_1$=（11,13,15,17,19）, 则 $S_1$的中位数是 15, 两 个序列的中位
        数是含它们所有元素的升序序列的中位数。例如，若 $S_2$ =（2,4,6,8,20）, 则$S_1$和$S_2$的中
        位数是 11。现在有两 个等长升序序列$A$和$B$, 试设计一个在时间和空间两 方面都尽可能
        高效的算法，找出两个序列 $A$和$B$的中位数。要求：
        \begin{subqitems}
            \subqitem 给出算法的基本设计思想
            \subqitem  根据设计思想，采用 C 或 C++或 Java 语言描述算法，关键之处给出注释
            \subqitem 说明你所设计算法的时间复杂度和空间复杂度
        \end{subqitems}
    \end{bbox}

\end{qitems}

\end{document}